% ============================================================
%  CLASE DEL DOCUMENTO
% ============================================================
\documentclass[12pt,oneside]{book}

% ============================================================
%  METADATOS
% ============================================================
\newcommand{\universidad}{Universidad de Buenos Aires}
\newcommand{\facultad}{Facultad de Derecho}
\newcommand{\carrera}{Abogacía / Traductorado Público}
\newcommand{\materia}{Inglés Jurídico --- Lectocomprensión de Textos Jurídicos en Inglés}
\newcommand{\titulo}{Apuntes de clase --- 22 de septiembre}
\newcommand{\autor}{Isaac Edgar Camacho Ocampo}
\newcommand{\anio}{2025}

% ============================================================
%  PAQUETES
% ============================================================
\usepackage[utf8]{inputenc}
\usepackage[T1]{fontenc}
\usepackage[spanish,es-nodecimaldot]{babel}

\usepackage[
    top=2.5cm,
    bottom=2.5cm,
    left=3cm,
    right=2.5cm,
    headheight=15pt
]{geometry}

\usepackage{amsmath}
\usepackage{amssymb}
\usepackage{mathtools}

\usepackage{graphicx}
\usepackage{float}
\usepackage{caption}
\usepackage{subcaption}

\usepackage{booktabs}
\usepackage{array}
\usepackage{longtable}
\usepackage{tabularx}

\usepackage{enumitem}

\usepackage{xcolor}
\usepackage[most]{tcolorbox}

\usepackage{fancyhdr}

\usepackage{ragged2e}

\usepackage[
    colorlinks=true,
    linkcolor=blue!50!black,
    citecolor=green!40!black,
    urlcolor=blue!60!black,
    pdftitle={Apuntes de Inglés Jurídico},
    pdfauthor={Isaac Edgar Camacho Ocampo},
    pdfsubject={Apuntes universitarios},
    bookmarks=true
]{hyperref}

% ============================================================
%  RUTAS DE IMÁGENES
% ============================================================
\graphicspath{
    {./img/}
    {./graficos/}
    {./figuras/}
}

% ============================================================
%  CONFIGURACIÓN TIPOGRÁFICA
% ============================================================
\setlength{\parindent}{0pt}
\setlength{\parskip}{0.5em}

\setcounter{tocdepth}{2}

\clubpenalty=10000
\widowpenalty=10000

\setlist[itemize]{topsep=0.3em, itemsep=0.2em}
\setlist[enumerate]{topsep=0.3em, itemsep=0.2em}

% ============================================================
%  ENCABEZADOS Y PIES
% ============================================================
\pagestyle{fancy}
\fancyhf{}
\fancyhead[L]{\nouppercase{\leftmark}}
\fancyhead[R]{\small Inglés Jurídico}
\fancyfoot[C]{\thepage}
\renewcommand{\headrulewidth}{0.4pt}
\renewcommand{\footrulewidth}{0pt}

\fancypagestyle{plain}{
    \fancyhf{}
    \fancyfoot[C]{\thepage}
    \renewcommand{\headrulewidth}{0pt}
}

% ============================================================
%  COMANDOS MATEMÁTICOS
% ============================================================
\newcommand{\abs}[1]{\left|#1\right|}
\newcommand{\norm}[1]{\left\lVert#1\right\rVert}
\newcommand{\vect}[1]{\mathbf{#1}}
\newcommand{\deriv}[2]{\frac{d#1}{d#2}}
\newcommand{\pderiv}[2]{\frac{\partial#1}{\partial#2}}

% ============================================================
%  CAJAS ACADÉMICAS
% ============================================================

% Definición
\newtcolorbox{definition}[1]{
    enhanced,
    breakable,
    title={Definición: #1},
    fonttitle=\bfseries,
    colback=blue!3,
    colframe=blue!50!black,
    boxrule=0.8pt,
    arc=2mm
}

% Importante
\newtcolorbox{important}{
    enhanced,
    breakable,
    title={Importante},
    fonttitle=\bfseries,
    colback=yellow!8,
    colframe=orange!70!black,
    boxrule=0.8pt,
    arc=2mm
}

% Ejemplo
\newtcolorbox{example}[1]{
    enhanced,
    breakable,
    title={Ejemplo: #1},
    fonttitle=\bfseries,
    colback=green!4,
    colframe=green!40!black,
    boxrule=0.8pt,
    arc=2mm
}

% Fórmula / Regla
\newtcolorbox{formula}[1]{
    enhanced,
    breakable,
    title={Fórmula / Regla: #1},
    fonttitle=\bfseries,
    colback=purple!4,
    colframe=purple!50!black,
    boxrule=0.8pt,
    arc=2mm
}

% Ejercicio
\newtcolorbox{exercise}[1]{
    enhanced,
    breakable,
    title={Ejercicio: #1},
    fonttitle=\bfseries,
    colback=gray!5,
    colframe=gray!60!black,
    boxrule=0.8pt,
    arc=2mm
}

% Nota
\newtcolorbox{note}{
    enhanced,
    breakable,
    title={Nota},
    fonttitle=\bfseries,
    colback=white,
    colframe=gray!50,
    boxrule=0.6pt,
    arc=2mm
}

% Fragmento trabajado en clase (con borde dorado, tal como en el original)
\newtcolorbox{fragment}[1]{
    enhanced,
    breakable,
    title={Fragmento trabajado en clase: #1},
    fonttitle=\bfseries,
    colback=yellow!4,
    colframe=yellow!60!black,
    boxrule=0.8pt,
    arc=2mm
}

% ============================================================
%  DOCUMENTO
% ============================================================
\begin{document}

% ============================================================
%  PORTADA
% ============================================================
\begin{titlepage}

\thispagestyle{empty}

\begin{center}

\vspace*{1cm}

{\large \textsc{\universidad}}

\vspace{0.2cm}

{\large \textsc{\facultad}}

\vspace{0.2cm}

{\large \carrera}

\vspace{3cm}

{\Huge \textbf{Inglés Jurídico}}

\vspace{0.5cm}

{\Large Lectocomprensión de Textos Jurídicos en Inglés}

\vspace{1cm}

{\LARGE \titulo}

\vspace{0.8cm}

\textit{Leer un texto jurídico en otra lengua es, ante todo, aprender a pensar en otro sistema.}

\vspace{1.2cm}

\rule{0.70\textwidth}{0.5pt}

\vspace{0.5cm}

{\large Apuntes de estudio}

\vspace{0.5cm}

\rule{0.70\textwidth}{0.5pt}

\vfill

{\large \autor}

\vspace{0.3cm}

{\large \anio}

\end{center}

\vfill

\begin{flushleft}
\textit{Este es un modesto aporte a los estudiantes de ``Inglés Jurídico'' no pretende sustituir el material oficial de la materia.}
\end{flushleft}

\end{titlepage}

% ============================================================
%  SECCIÓN DE ORIENTACIÓN INICIAL
% ============================================================
\frontmatter

\chapter*{Relevancia profesional de este material}
\addcontentsline{toc}{chapter}{Relevancia profesional de este material}

Quien traduce o litiga en un entorno bilingüe necesita algo más que vocabulario: necesita entender que un mismo término puede significar cosas distintas en dos sistemas jurídicos. Confundir \textit{criminal} con culpable, traducir \textit{property} como inmueble, o asimilar \textit{negligence} a la categoría de culpa del derecho continental sin los matices correspondientes puede cambiar el sentido de un contrato, una demanda o un dictamen.

Esta clase entrena precisamente esa competencia: leer un texto jurídico en inglés sin traducir palabra por palabra, reconociendo la lógica del sistema del \textit{common law} ---sus roles procesales, sus remedios, sus categorías de responsabilidad--- para poder trasladarla correctamente al derecho local. Es una herramienta directa para la traducción pública, el ejercicio profesional con clientes o estudios extranjeros, la redacción de contratos internacionales y la lectura de jurisprudencia comparada.

Más allá de lo estrictamente profesional, este tipo de lectura cultiva un hábito valioso: no asumir que las palabras significan lo mismo en todos los contextos, y verificar el sentido antes de actuar sobre él.

\chapter*{Introducción: un mismo viaje, cuatro etapas}
\addcontentsline{toc}{chapter}{Introducción: un mismo viaje, cuatro etapas}

El contenido de esta clase puede pensarse como el mapa de un viaje judicial. En primer lugar, es necesario saber qué tipo de viaje es el que se emprende: no es lo mismo un proceso penal, donde el Estado acciona contra una persona, que un proceso civil, entre particulares; y cada uno conduce a un destino distinto ---una pena o una reparación--- (Bloque~1). Dentro del terreno civil aparece luego un tipo de conflicto frecuente y propio: el daño entre particulares, denominado \textit{tort}, con sus reglas específicas de responsabilidad (Bloque~2). Para recorrer ese camino se requiere un guía: el abogado, que en inglés adopta múltiples denominaciones según el rol que cumple (Bloque~3). Finalmente, se trabaja el itinerario concreto: los primeros pasos formales del proceso, es decir, cómo se inicia una demanda civil (Bloque~4). Los cuatro bloques son, en definitiva, las piezas de un mismo recorrido: del tipo de proceso, a su conflicto típico, a quien lo conduce, a cómo se pone en marcha.

% ============================================================
%  ÍNDICE
% ============================================================
\tableofcontents

% ============================================================
%  CUERPO PRINCIPAL
% ============================================================
\mainmatter

% ============================================================
%  CAPÍTULO 1
% ============================================================
\chapter{Fundamentos del derecho anglosajón: proceso civil y proceso penal}

\section{El texto de trabajo y la actividad de anticipación}

La clase trabaja sobre el texto titulado \textit{Steps in a Civil Lawsuit}, dedicado a los pasos del proceso civil en el sistema del \textit{common law}. Antes de abordar el cuerpo del texto, se propone una actividad de anticipación a partir de la \textbf{macroestructura}: título y subtítulos.

Los subtítulos identificados en el texto son: \textit{complaint and answer}, \textit{motions in the early stages} y \textit{discovery and settlements}. Dado el alto grado de tecnicismo del tema, se advierte que la sola lectura de esos elementos paratextuales puede no resultar suficientemente orientadora para quien desconoce el sistema. Por ello, la lectura inicial se limita al primer segmento del texto (\textit{initial steps}), que funciona como introducción general.

\section{La distinción entre causa penal y causa civil}

El texto parte de una observación sobre la influencia de las series de televisión y la cobertura periodística de juicios en la percepción pública del derecho: como consecuencia de esa exposición mediática, el público está más familiarizado con el proceso penal que con el civil. A partir de esa constatación, el texto traza la distinción fundamental entre ambos tipos de proceso.

\begin{important}
En inglés jurídico, la expresión \textbf{\textit{criminal case}} no debe traducirse pensando en la noción de ``criminal'' en español. En castellano, el término remite a la idea de culpabilidad. En inglés jurídico, en cambio, \textit{criminal} es un adjetivo neutro que simplemente indica la rama del derecho involucrada, equivalente a ``penal''. Para evitar confusiones en la traducción, corresponde utilizar ``penal'' y reservar ``criminal'' para contextos donde el sistema de origen así lo exija.
\end{important}

\subsection{Las partes del proceso}

En una \textbf{causa penal} (\textit{criminal case}), el proceso es impulsado por el Estado a través del \textbf{fiscal} (\textit{prosecutor}), frente al \textbf{acusado o imputado} (\textit{defendant}), cuya denominación varía según el momento procesal.

En una \textbf{causa civil} (\textit{civil case}), intervienen dos particulares: el \textbf{demandante} (\textit{plaintiff}) y el \textbf{demandado} (\textit{defendant}). El término \textit{defendant} se mantiene en ambas ramas; quien cambia de denominación es la parte contraria: \textit{prosecutor} en lo penal, \textit{plaintiff} en lo civil.

\section{Las consecuencias de cada proceso: penas y reparaciones}

Cada rama del derecho produce consecuencias distintas:

\textbf{En el ámbito penal}, la consecuencia es una pena (\textit{penalty} o \textit{punishment}). Las principales formas son:
\begin{itemize}
    \item \textit{Imprisonment} / \textit{incarceration}: pena privativa de libertad.
    \item \textit{Fine}: multa dineraria.
    \item \textit{Capital punishment}: pena de muerte, en los estados donde está contemplada.
    \item \textit{Probation} y tareas comunitarias: figuras que también implican una cierta pérdida o restricción de libertad.
\end{itemize}

\textbf{En el ámbito civil}, la reparación por excelencia frente al daño (\textit{injury}) es \textit{damages}, es decir, una indemnización dineraria por daños y perjuicios. Existen además otras dos formas de reparación posibles:
\begin{itemize}
    \item \textit{Specific performance}: cumplimiento específico de una obligación; equivale, por ejemplo, a un derecho a réplica cuando la reparación dineraria no resulta suficiente o adecuada.
    \item \textit{Injunction}: medida de no hacer; ordena el cese de una determinada conducta.
\end{itemize}

El derecho civil, tal como lo organiza el texto, comprende distintas materias: cuestiones de familia (\textit{family matters}), cuestiones contractuales (\textit{contract law}) y el denominado \textit{tort law}, desarrollado en el capítulo siguiente.

\bigskip

\begin{table}[H]
\centering
\caption{Comparación entre proceso penal y proceso civil en el \textit{common law}}
\begin{tabularx}{\textwidth}{
    >{\raggedright\arraybackslash}p{0.22\textwidth}
    >{\raggedright\arraybackslash}X
    >{\raggedright\arraybackslash}X
}
\toprule
\textbf{Criterio} & \textbf{Proceso penal (\textit{Criminal case})} & \textbf{Proceso civil (\textit{Civil case})} \\
\midrule
Partes & Estado (a través del \textit{prosecutor}) vs. \textit{defendant} & \textit{Plaintiff} (actor) vs. \textit{defendant}; dos particulares \\
\addlinespace
Naturaleza del conflicto & Delito contra el orden público / la sociedad & Disputa privada por un daño o agravio (\textit{injury}) sufrido por una de las partes \\
\addlinespace
Consecuencia típica & \textit{Penalty} / \textit{punishment}: \textit{imprisonment}, \textit{fine}, \textit{capital punishment}, \textit{probation} & \textit{Damages} (indemnización), \textit{specific performance}, \textit{injunction} \\
\addlinespace
Materias comprendidas & Delitos tipificados por la ley penal & \textit{Family law}, \textit{contract law}, \textit{tort law} \\
\bottomrule
\end{tabularx}
\end{table}

\section{Las finalidades de la pena}

El texto enuncia cuatro finalidades de la pena que el derecho penal anglosajón reconoce: \textit{retribution}, \textit{deterrence}, \textit{incapacitation} y \textit{rehabilitation}.

\begin{important}
El término \textbf{\textit{retribution}} puede generar una asociación equivocada con la idea coloquial de venganza (``ojo por diente''). En el lenguaje jurídico, no remite a ese sentido personal, sino a una lógica redistributiva: quien causó un daño debe repararlo. Es llamativo que el lenguaje jurídico recurra a un término que en el habla cotidiana evoca la venganza para nombrar esta finalidad.
\end{important}

\begin{table}[H]
\centering
\caption{Las cuatro finalidades de la pena en el derecho anglosajón}
\begin{tabularx}{\textwidth}{
    >{\raggedright\arraybackslash}p{0.26\textwidth}
    >{\raggedright\arraybackslash}X
    >{\raggedright\arraybackslash}p{0.28\textwidth}
}
\toprule
\textbf{Finalidad} & \textbf{Explicación} & \textbf{Ejemplo o aclaración} \\
\midrule
\textit{Retribution} (retribución) & No equivale a venganza personal, sino a una lógica redistributiva: quien causó un daño debe repararlo. & Se compara y distingue del sentido coloquial de ``ojo por diente''. \\
\addlinespace
\textit{Deterrence} (disuasión) & \textit{To deter someone from doing something} es disuadirlo. Busca que ni el infractor ni la sociedad repitan la conducta; funciona como efecto ejemplificador. & ``No paso el semáforo en rojo porque puedo chocar y, además, porque me van a poner una multa'': la sanción conocida disuade la conducta. \\
\addlinespace
\textit{Incapacitation} (incapacitación, protección de la comunidad) & Apunta a impedir que la persona vuelva a dañar a la sociedad, apartándola de la posibilidad de reiterar la conducta. & Quien ejerce la medicina sin título es inhabilitado; un policía es apartado de su cargo. \\
\addlinespace
\textit{Rehabilitation} (rehabilitación) & Tiene un fin social: la reinserción de la persona en la sociedad. & Equivale a lo que en español se denomina ``reinserción''. \\
\bottomrule
\end{tabularx}
\end{table}

\begin{example}{Incapacitación en la práctica}
A propósito de la \textit{incapacitation}, puede ilustrarse el concepto con el caso de empleados de un banco que, aprovechando el acceso a las cajas de seguridad de los clientes, sustraían dinero en pequeñas sumas. La sanción consistente en apartar a esas personas de su función cumple una finalidad de incapacitación: impide que la conducta se repita al retirar a quienes la ejecutaban del lugar que se los permitía.
\end{example}

\begin{note}
El análisis de las finalidades de la pena es de carácter más doctrinario. El desarrollo pormenorizado de los tipos de pena en particular se retomará al avanzar sobre la materia penal.
\end{note}

% ============================================================
%  CAPÍTULO 2
% ============================================================
\chapter{Derecho de daños: \textit{Tort Law}}

\section{Primer acercamiento al término}

Al mencionar el texto que el derecho civil comprende el \textit{tort law} como una de sus ramas (junto con \textit{family law} y \textit{contract law}), corresponde advertir que la sola forma de la palabra \textit{tort} puede generar asociaciones erróneas en español, por su parecido fonético con ``torta'' o ``tortura''. Ninguna de esas asociaciones es pertinente.

\section{Definición de \textit{tort}}

\begin{definition}{Tort}
\textit{``A tort is a wrongful act or omission, other than a breach of contract or trust, that injures another person, their property or their reputation, entitling the injured party to compensation.''}

\medskip

Un \textit{tort} es un acto u omisión ilícita ---distinta del incumplimiento de un contrato o de un abuso de confianza--- que causa un daño a otra persona, a sus bienes o a su reputación, y que da derecho a la parte damnificada a una compensación.
\end{definition}

\section{Análisis de los elementos de la definición}

\subsection{\textit{Wrong} / \textit{wrongful}: lo antijurídico, no lo inmoral}

En el lenguaje jurídico, \textbf{\textit{wrong}} refiere a algo ilícito, antijurídico, contrario a la norma o contrario a la ley. La idea de \textit{wrong} en este contexto va más allá de la noción moral o ética de ``algo que está mal'': designa lo que es contrario al derecho. En ocasiones se lo utiliza incluso como sinónimo de ilícito o delito.

El acto ilícito que constituye un \textit{tort} queda expresamente excluido del incumplimiento contractual (\textit{breach of contract}) y del abuso de confianza (\textit{breach of trust}), entendido este último como el quebrantamiento de la confianza propia de un vínculo, un cargo o una jerarquía entre las personas.

\subsection{\textit{Injury}: daño en sentido amplio}

\begin{important}
\textbf{\textit{Injury}} es un término que puede inducir a error. En el lenguaje cotidiano se asocia a la idea de una lastimadura o lesión física. En el lenguaje jurídico es bastante más amplio: designa un agravio, una lesión o un perjuicio que \textbf{no necesariamente tiene que ser físico}. Se trata, por lo tanto, de un falso amigo parcial que requiere atención en la traducción.
\end{important}

\subsection{\textit{Property}: patrimonio, no solamente inmuebles}

\begin{important}
\textbf{\textit{Property}} es otro término que puede inducir a error. En español, ``propiedad'' evoca principalmente la idea de un inmueble. En el contexto de la definición de \textit{tort}, \textit{property} refiere en términos generales al patrimonio de una persona: abarca tanto bienes muebles como inmuebles.
\end{important}

\subsection{\textit{Entitled to compensation}: el derecho a ser reparado}

El último elemento de la definición ---\textit{entitling the injured party to compensation}--- establece que el daño causado habilita a la parte damnificada a reclamar una compensación. Este derecho no opera de manera automática: es necesario analizar la conducta de la propia víctima y su eventual incidencia en el resultado dañoso.

\begin{example}{Accidente de tránsito y responsabilidad compartida}
Si la persona que sufrió el daño fue de alguna manera partícipe del resultado ---por ejemplo, si no realizó correctamente una maniobra de giro y parte del accidente guarda relación con su propia conducta---, el derecho a la compensación (\textit{entitled to compensation}) exige analizar si a esa persona le corresponde efectivamente ese derecho, o si el hecho de haberse puesto en situación de riesgo incide en la reparación que puede reclamar.
\end{example}

El mismo criterio se aplica a la responsabilidad del fabricante frente a un producto defectuoso ---por ejemplo, un artefacto eléctrico que causa un daño por su mal funcionamiento---, materia que el texto aborda bajo la categoría de \textit{product liability}.

\section{Clasificación de los torts: dolosos y culposos}

El texto clasifica los \textit{torts} en dos grandes categorías:

\begin{table}[H]
\centering
\caption{Clasificación de los \textit{torts}}
\begin{tabularx}{\textwidth}{
    >{\raggedright\arraybackslash}p{0.26\textwidth}
    >{\raggedright\arraybackslash}X
    >{\raggedright\arraybackslash}p{0.30\textwidth}
}
\toprule
\textbf{Clasificación} & \textbf{Definición} & \textbf{Ejemplos del texto} \\
\midrule
\textit{Intentional torts} (dolosos) & Existe intención de causar el daño. & \textit{Battery} (agresión física), \textit{assault} (amenaza de agresión), \textit{libel} (difamación escrita) y \textit{slander} (difamación oral). \\
\addlinespace
\textit{Negligent torts} (culposos) & El daño no fue querido: es resultado de una conducta que causa perjuicios no intencionales. & \textit{Auto accidents}, \textit{medical malpractice}, \textit{product liability}. \\
\bottomrule
\end{tabularx}
\end{table}

Dentro de los \textit{intentional torts}, conviene precisar la distinción entre \textbf{\textit{assault}} y \textbf{\textit{battery}}:
\begin{itemize}
    \item \textit{Assault}: amenaza de agresión; puede existir sin que llegue a producirse el golpe.
    \item \textit{Battery}: concreción efectiva de la agresión física; puede darse sin amenaza previa.
\end{itemize}

Respecto de \textit{libel} y \textit{slander}, ambos remiten a la difamación (calumnias e injurias), y su distinción depende del medio empleado: \textit{libel} es la difamación escrita y \textit{slander} la oral.

\section{\textit{Tort law} frente a las categorías del sistema continental: \textit{negligence} y culpa}

\begin{important}
\textbf{\textit{Negligence}} no equivale exactamente a la categoría de ``culpa'' del derecho continental. El \textit{common law} utiliza \textit{negligence} como una categoría amplia que abarca todo lo que no es intencional, sin efectuar la subdivisión más fina que existe en el derecho local entre \textbf{negligencia}, \textbf{imprudencia} e \textbf{impericia}. Intentar analizar la \textit{negligence} del sistema anglosajón con los mismos matices que el sistema continental puede conducir a errores de interpretación.
\end{important}

La diferencia puede ilustrarse con el caso de la mala praxis médica: en el derecho anglosajón, el análisis se reduce a establecer si existió intención de causar el daño (lo que lo ubicaría entre los \textit{intentional torts}) o si el daño fue consecuencia del deficiente ejercicio de una técnica, sin intención de dañar (lo que lo coloca dentro de los \textit{negligent torts}). No se realizan subdivisiones ulteriores entre negligencia, imprudencia e impericia.

\begin{note}
A diferencia del derecho continental, donde estas cuestiones se integran dentro del derecho civil en general, en el mundo anglosajón ---y especialmente en los Estados Unidos--- el \textit{tort law} llegó a constituirse en una rama bastante autónoma y elaborada del derecho, que da lugar a demandas de gran resonancia pública (como las dirigidas contra cadenas de comida rápida). Estas acciones no configuran delitos penales ni incumplimientos contractuales, sino ilícitos civiles propios del \textit{tort law}.
\end{note}

El texto recomienda consultar a un abogado especializado en el área específica del conflicto (\textit{attorney particular}): uno especializado en derecho penal para un proceso penal, uno de familia para una cuestión familiar, etcétera. Cada rama tiene sus propias reglas, aunque los principios básicos del derecho se apliquen transversalmente.

% ============================================================
%  CAPÍTULO 3
% ============================================================
\chapter{Vocabulario de la profesión jurídica en inglés}

\section{Los términos generales para nombrar al abogado}

El inglés dispone de varias voces para referirse a quien ejerce la abogacía. Si bien no son exactamente sinónimas, a grandes rasgos remiten a la misma figura, con matices de uso según el contexto y el rol que el profesional cumple en cada situación. Los principales términos son: \textit{lawyer}, \textit{attorney}, \textit{counsel}, \textit{barrister}, \textit{solicitor}, \textit{squire} (o \textit{esquire}) y \textit{advocate}.

\section{El sistema del Reino Unido: \textit{barrister} y \textit{solicitor}}

En el Reino Unido, quien egresa de la carrera de derecho enfrenta dos caminos posibles de ejercicio profesional, en función de si rinde o no el examen ante la \textbf{\textit{Bar Association}} ---el colegio profesional que habilita para litigar.

\begin{definition}{Barrister}
Abogado que aprobó el examen de la \textit{Bar Association} y, en consecuencia, está habilitado para \textbf{litigar y comparecer ante los tribunales}, patrocinando la causa. Es análogo a la figura del abogado patrocinante en el derecho local.
\end{definition}

\begin{definition}{Solicitor}
Abogado que puede \textbf{asesorar al cliente}, reunirse con él, redactar escritos e investigar, pero \textbf{no puede comparecer ante el tribunal}. Lleva adelante gran parte del trabajo preparatorio de la causa. Es análogo, en ciertos aspectos, a la figura del procurador en el derecho local.
\end{definition}

\begin{important}
Ambos ---\textit{barrister} y \textit{solicitor}--- son abogados. La diferencia no radica en el título sino en la \textbf{habilitación para litigar}. La situación es parcialmente comparable a la de aquellos abogados que, en el sistema local, no tramitan su matrícula porque se dedican a tareas de asesoramiento societario y no a la litigación.
\end{important}

\begin{table}[H]
\centering
\caption{Comparación entre \textit{barrister} y \textit{solicitor} en el sistema del Reino Unido}
\begin{tabularx}{\textwidth}{
    >{\raggedright\arraybackslash}p{0.36\textwidth}
    >{\raggedright\arraybackslash}X
    >{\raggedright\arraybackslash}X
}
\toprule
\textbf{Aspecto} & \textbf{\textit{Barrister}} & \textbf{\textit{Solicitor}} \\
\midrule
Litigar / comparecer ante tribunales & Sí & No \\
\addlinespace
Asesorar y reunirse con el cliente & Puede hacerlo & Sí, es su función principal \\
\addlinespace
Redactar escritos, investigar & Puede hacerlo & Sí \\
\addlinespace
Habilitación & Aprobó el examen de la \textit{Bar Association} & No rindió, no aprobó, o eligió no litigar \\
\addlinespace
Figura análoga en el derecho local & Abogado patrocinante & Procurador \\
\bottomrule
\end{tabularx}
\end{table}

\section{\textit{Attorney}: \textit{at law} versus \textit{in fact}}

El término \textbf{\textit{attorney}} puede resultar ambiguo si no se especifica su tipo:

\begin{itemize}
    \item \textbf{\textit{Attorney at law}}: equivale a abogado en sentido pleno.
    \item \textbf{\textit{Attorney in fact}}: designa a un \textbf{apoderado}; una persona a quien se le otorgó un poder para representar a otra, sin que sea necesariamente abogado.
\end{itemize}

\begin{example}{El apoderado como \textit{attorney in fact}}
Una persona jubilada puede designar un apoderado (\textit{attorney in fact}) para realizar trámites bancarios en su nombre, sin que ese apoderado sea necesariamente abogado. También puede darse el caso de un abogado que ejerce simultáneamente como apoderado de su propio cliente, reuniendo así ambas condiciones.
\end{example}

\section{Otros términos según la función: \textit{counsel} y \textit{advocate}}

El término empleado para denominar al abogado varía en función del contexto y el rol que el profesional cumple:

\begin{table}[H]
\centering
\caption{Términos para denominar al abogado en inglés jurídico}
\begin{tabularx}{\textwidth}{
    >{\raggedright\arraybackslash}p{0.26\textwidth}
    >{\raggedright\arraybackslash}X
}
\toprule
\textbf{Término} & \textbf{Rol o contexto de uso} \\
\midrule
\textit{Lawyer} / \textit{Attorney at law} & Término general para abogado. \\
\addlinespace
\textit{Attorney in fact} & Apoderado; no necesariamente abogado. \\
\addlinespace
\textit{Counsel} & Abogado patrocinante en una demanda o instrumento judicial (ej.: \textit{counsel for plaintiff}). \\
\addlinespace
\textit{Barrister} & Abogado habilitado para litigar y comparecer ante el tribunal (Reino Unido). \\
\addlinespace
\textit{Solicitor} & Abogado que asesora, prepara y redacta, sin comparecer ante el tribunal (Reino Unido). \\
\addlinespace
\textit{Advocate} & Abogado que asesora, típicamente en contextos empresariales o de negocios, sin litigar. \\
\addlinespace
\textit{Squire} / \textit{Esquire} & Tratamiento formal que acompaña el nombre del abogado, análogo a ``doctor''. \\
\bottomrule
\end{tabularx}
\end{table}

% ============================================================
%  CAPÍTULO 4
% ============================================================
\chapter{El proceso civil: etapas iniciales de la demanda}

\section{El \textit{complaint} como primer acto procesal}

El proceso civil comienza con la \textbf{demanda}, denominada en inglés \textbf{\textit{complaint}}. Sin embargo, antes de que exista un \textit{complaint} debe haber previamente un conflicto o controversia (\textit{dispute}), no necesariamente mediada por un proceso de mediación formal.

\begin{definition}{\textit{Complaint}}
El \textit{complaint} es el \textbf{primer acto procesal} del proceso civil; el documento inicial que da comienzo a la acción (\textit{initial document initiating an action}). Es presentado por el \textit{plaintiff} o parte actora. Aunque en un diccionario no jurídico el verbo \textit{to complain} se traduce como ``quejarse'', en su uso técnico-procesal el \textit{complaint} equivale a la demanda.
\end{definition}

\section{Los elementos del \textit{complaint}: \textit{cause of action}, \textit{relief} y \textit{allegations}}

\subsection{\textit{Cause of action}: el fundamento de la acción}

\begin{definition}{\textit{Cause of action}}
El \textit{cause of action} es el concepto jurídico asimilable a lo que en el derecho local se denomina \textbf{pretensión}: es el fundamento por el cual se presenta la demanda. No alcanza con afirmar que se sufrió un daño; es necesario establecer el nexo entre los hechos y la persona demandada. Ese nexo debe apoyarse en la existencia de un deber de cuidado incumplido, en una obligación de hacer o de no hacer que la conducta del demandado desatendió, y en la correspondencia entre ese daño y esa conducta. Si la pretensión no está bien fundada, la demanda puede no prosperar; por eso los jueces analizan especialmente este punto al recibir una demanda.
\end{definition}

\subsection{\textit{Damages or relief}: el pedido de reparación}

Junto con la causa de la acción, el \textit{complaint} incluye el pedido concreto de reparación. El texto lo expresa así: \textit{the complaint asks for damages or relief from the defendant}. Se trata de la pretensión resarcitoria o del remedio que se reclama frente al daño sufrido.

\subsection{\textit{Allegations}: las imputaciones aún no probadas}

El \textit{complaint} identifica a la persona señalada como responsable del daño mediante el uso del verbo \textbf{\textit{to allege}} y sus derivados (\textit{the alleged conduct}, \textit{the alleged person responsible}).

\begin{important}
\textbf{\textit{To allege}} implica afirmar algo que aún no ha sido corroborado. Refleja la \textbf{presunción de inocencia}: hasta que se demuestre lo contrario, la responsabilidad es una imputación, no una certeza establecida. Esta distinción tiene consecuencias prácticas: un periodista que afirme la responsabilidad de una persona sin emplear este resguardo lingüístico puede enfrentarse a consecuencias legales.
\end{important}

\subsection{Contenido fáctico y legal}

El \textit{complaint} incluye además una narración de cuestiones tanto fácticas (\textit{factual}) como legales (\textit{legal}). No basta con relatar los hechos: es necesario también encuadrarlos normativamente.

\begin{example}{El caso del perro}
Una persona es mordida por un perro de una raza considerada peligrosa, que su dueño paseaba sin correa ni sujeción alguna. La demanda deberá narrar los hechos (el ataque, las circunstancias) y, además, el fundamento legal: por ejemplo, que se trata de un animal de raza peligrosa, o que existía una obligación normativa de mantenerlo con correa que fue incumplida.

Variante para analizar el nexo causal: si fuera la propia víctima quien, estando el perro detrás de una reja, introdujo la mano dentro de la propiedad, corresponde analizar hasta qué punto el daño resulta atribuible a la conducta del dueño o, en cambio, a la propia conducta de la víctima.
\end{example}

\section{Los \textit{pleadings} y los actos procesales que siguen al \textit{complaint}}

El \textit{complaint} constituye el primer escrito del proceso (\textit{first pleading}). Los \textit{pleadings} en plural, sin embargo, no se agotan en ese primer acto. A la demanda le siguen, entre otros:

\begin{itemize}
    \item La notificación.
    \item El traslado.
    \item La contestación de demanda.
    \item Eventualmente, la reconvención o contrademanda.
\end{itemize}

Estos contenidos se anuncian como materia de la clase siguiente.

% ============================================================
%  CAPÍTULO 5: CUADRO CORNELL
% ============================================================
\chapter{Cuadro de estudio Cornell}

El siguiente cuadro organiza el contenido de los cuatro bloques temáticos en el formato Cornell: preguntas o conceptos disparadores en la columna izquierda, y el desarrollo correspondiente en la columna derecha. La última fila contiene la síntesis integradora de todo lo trabajado.

\begin{table}[H]
\centering
\caption{Cuadro Cornell --- Inglés Jurídico, clase del 22 de septiembre}
\begin{tabularx}{\textwidth}{
    >{\raggedright\arraybackslash}p{0.30\textwidth}
    >{\raggedright\arraybackslash}X
}
\toprule
\textbf{Preguntas / palabras clave} & \textbf{Notas / respuestas} \\
\midrule

\textbf{¿Por qué no debe traducirse \textit{criminal case} pensando en culpabilidad?} &
En inglés jurídico, \textit{criminal} es un adjetivo neutro que solo indica la rama del derecho (equivalente a ``penal''). No implica, como en español, que la persona sea culpable. \\
\addlinespace

\textbf{¿Cuáles son las partes de un proceso penal y de un proceso civil?} &
En el proceso penal: el Estado (a través del \textit{prosecutor}) contra el \textit{defendant}. En el proceso civil: dos particulares, el \textit{plaintiff} contra el \textit{defendant}. \\
\addlinespace

\textbf{¿Qué consecuencias corresponden a cada rama?} &
Penal: \textit{penalty} / \textit{punishment} (\textit{imprisonment}, \textit{fine}, \textit{capital punishment}, \textit{probation}). Civil: \textit{damages}, \textit{specific performance}, \textit{injunction}. \\
\addlinespace

\textbf{¿Cuáles son las cuatro finalidades de la pena?} &
\textit{Retribution} (lógica redistributiva, no venganza), \textit{deterrence} (disuasión, efecto ejemplificador), \textit{incapacitation} (impedir que la persona repita el daño; protección de la comunidad) y \textit{rehabilitation} (reinserción social). \\
\addlinespace

\textbf{¿Qué es un \textit{tort}?} &
Un acto u omisión ilícita ---distinta del incumplimiento contractual o del abuso de confianza--- que causa un daño a otra persona, a sus bienes o a su reputación, y que da derecho a compensación. \\
\addlinespace

\textbf{¿Por qué \textit{injury} y \textit{property} son falsos amigos en este contexto?} &
\textit{Injury} no designa solo una lesión física, sino cualquier agravio o perjuicio. \textit{Property} no refiere solo a inmuebles, sino al patrimonio o los bienes muebles e inmuebles de una persona. \\
\addlinespace

\textbf{¿Cómo se clasifican los \textit{torts}?} &
En \textit{intentional torts} (dolosos: \textit{battery}, \textit{assault}, \textit{libel}, \textit{slander}) y \textit{negligent torts} (culposos: \textit{auto accidents}, \textit{medical malpractice}, \textit{product liability}). \\
\addlinespace

\textbf{¿Por qué \textit{negligence} no equivale exactamente a la ``culpa'' del derecho continental?} &
El \textit{common law} utiliza \textit{negligence} como categoría amplia para todo lo no doloso, sin la subdivisión más fina que existe en el derecho local entre negligencia, imprudencia e impericia. \\
\addlinespace

\textbf{¿Cuál es la diferencia entre \textit{barrister} y \textit{solicitor}?} &
El \textit{barrister} rindió el examen de la \textit{Bar Association} y puede litigar y comparecer ante el tribunal. El \textit{solicitor} asesora, prepara y redacta, pero no comparece; es análogo al procurador. \\
\addlinespace

\textbf{¿Qué diferencia a un \textit{attorney at law} de un \textit{attorney in fact}?} &
El \textit{attorney at law} es abogado en sentido propio. El \textit{attorney in fact} es un apoderado: una persona con poder para representar a otra, sin ser necesariamente abogado. \\
\addlinespace

\textbf{¿Qué elementos debe contener un \textit{complaint}?} &
El \textit{cause of action} (fundamento y nexo con el demandado), el pedido de reparación (\textit{damages or relief}), las \textit{allegations} (imputaciones aún no probadas) y la narración de cuestiones \textit{legal and factual}. \\
\addlinespace

\textbf{¿Qué actos procesales siguen al \textit{complaint}?} &
La notificación, el traslado, la contestación de demanda y, eventualmente, la reconvención o contrademanda. \\

\midrule

\multicolumn{2}{>{\arraybackslash}p{\textwidth}}{%
\textbf{Síntesis:} La clase reconstruyó la lógica del \textit{common law} para tramitar un conflicto, partiendo de la distinción fundamental entre proceso penal y proceso civil ---diferenciados en sus partes, su propósito y sus consecuencias--- y descendiendo progresivamente hacia el detalle. Se vio que la pena combina cuatro finalidades: retribución, disuasión, incapacitación y rehabilitación. Dentro del universo civil se aisló el \textit{tort} como categoría propia del daño entre particulares, distinta del incumplimiento contractual, y se clasificó en dolosos y culposos, subrayando que \textit{negligence} no coincide exactamente con la ``culpa'' del derecho local. Se repasó la familia de términos que el inglés emplea para nombrar al abogado (\textit{barrister}, \textit{solicitor}, \textit{attorney}, \textit{counsel}, \textit{advocate}), y se trabajaron los elementos esenciales del primer acto procesal civil: el \textit{complaint}, con su \textit{cause of action}, su pedido de reparación y su narración fáctica y legal. El hilo conductor fue una advertencia constante: las palabras que más se parecen al español (\textit{criminal}, \textit{injury}, \textit{property}, \textit{negligence}) son, paradójicamente, las que exigen mayor cuidado, porque el sistema jurídico que las produjo no organiza siempre los conceptos de la misma manera que el propio.%
} \\

\bottomrule
\end{tabularx}
\end{table}

\end{document}
